% =====================================================================
%  ANEXO B. Especificación de procedimientos de prueba
% =====================================================================
\section{Especificación de procedimientos de prueba}
\label{anx:procedimientos}

\begin{notanormativa}[Diferencia con el caso de prueba]
Un caso de prueba especifica una verificación concreta. Un \textbf{procedimiento} agrupa y
\emph{ordena} varios casos en una secuencia ejecutable, añadiendo la preparación previa y la
limpieza posterior. Corresponde a la actividad \textbf{TD6} de \normap{2}: es el artefacto que
realmente se ejecuta.
\end{notanormativa}

\begin{observacion}[Anexo sin formato]
\obsproblema{El anexo equivalente de la plantilla anterior pedía «incluir los procedimientos de
prueba aplicados» sin proporcionar campos ni estructura.}
\obscambio{Ficha con identificador, propósito, casos incluidos y justificación de su orden,
entorno y datos requeridos, preparación, pasos, guion asociado y limpieza, más un índice.}
\obsfuente{OBS §2 y §6 · PDF §1.4 (TD6)}
\end{observacion}

\begin{instruccion}
Reproduzca la ficha siguiente por cada procedimiento. El orden de los casos importa: indique si
existen dependencias que obliguen a una secuencia concreta. La preparación y la limpieza son las
que hacen repetible el procedimiento; si se omiten, la segunda ejecución parte de un estado
distinto de la primera y los veredictos no son comparables.
\end{instruccion}

\subsection*{Ficha de procedimiento de prueba}

\begin{xltabular}{\textwidth}{|E{4.4cm}|Y|}
\caption{Formato de especificación de un procedimiento de prueba.}\label{tab:ficha-procedimiento}\\ \hline
\endfirsthead
\multicolumn{2}{@{}l@{}}{\footnotesize\itshape\tablename~\thetable{} (continuación)}\\ \hline
\endhead
\multicolumn{2}{@{}r@{}}{\footnotesize\itshape continúa en la página siguiente}\\
\endfoot
\endlastfoot
  Identificador del procedimiento & \campo{PP-001} \\ \hline
  Versión                         & \campo{v1} \\ \hline
  Propósito                       & \campo{Qué conjunto de verificaciones agrupa y con qué fin} \\ \hline
  Nivel de prueba                 & \campo{Componente / integración / sistema / aceptación} \\ \hline
  Casos incluidos y orden         & \campo{1.~CP-001 \quad 2.~CP-004 \quad 3.~CP-002} \\ \hline
  Justificación del orden         & \campo{Dependencias entre casos o restricciones del entorno} \\ \hline
  Entorno requerido               & \campo{ENT-xx (Cuadro~\ref{tab:req-entorno})} \\ \hline
  Datos requeridos                & \campo{DAT-xx (Cuadro~\ref{tab:datos})} \\ \hline
  Preparación previa              & \campo{Pasos para dejar el sistema en el estado inicial} \\ \hline
  Pasos de ejecución              & \campo{Secuencia detallada, manual o automatizada} \\ \hline
  Guion asociado                  & \campo{Ruta e identificación de versión del script, si aplica} \\ \hline
  Limpieza y restauración         & \campo{Qué se deshace al terminar para poder repetir} \\ \hline
  Duración estimada               & \campo{minutos/horas} \\ \hline
  Autor y fecha                   & \campo{Nombre — dd/mm/aaaa} \\ \hline
\end{xltabular}

El Cuadro~\ref{tab:ficha-procedimiento} define la unidad que se planifica y se ejecuta; su
identificador es el que aparece en el registro de ejecución del \ref{anx:ejecuciones} y en el
cronograma del apartado~\ref{sec:calendario}.

\subsection*{Índice de procedimientos}

\begin{longtable}{|C{1.8cm}|L{4.6cm}|L{4.0cm}|C{2.0cm}|}
\caption{Índice de procedimientos de prueba especificados.}
\label{tab:indice-procedimientos} \\
\hline
\filaenc \textbf{ID} & \textbf{Propósito} & \textbf{Casos incluidos} & \textbf{Nivel} \\ \hline
\endfirsthead
\hline
\filaenc \textbf{ID} & \textbf{Propósito} & \textbf{Casos incluidos} & \textbf{Nivel} \\ \hline
\endhead
\hline \multicolumn{4}{|r|}{\footnotesize\textit{continúa en la página siguiente}} \\ \hline
\endfoot
\hline
\endlastfoot
\campo{PP-001} & \campo{Propósito} & \campo{CP-001, CP-004} & \campo{Sistema} \\ \hline
 & & & \\ \hline
 & & & \\ \hline
\end{longtable}

El Cuadro~\ref{tab:indice-procedimientos} permite comprobar que todos los casos especificados
en el Cuadro~\ref{tab:indice-casos} están asignados a algún procedimiento; un caso huérfano
nunca llegará a ejecutarse.
